% ======================================================================
% SECTION 2: METHODS
% Final prose only. No bracketed instructions in this section, per the
% project brief. The next Claude Code 4.7 Max pass should leave this
% section untouched except to fix typographical errors flagged by the
% global formatting brief in main.tex.
% ======================================================================

\subsection{AI Generations and Author Roles}
\label{sec:methods-ai}

Primary AI generations to create the first relevant paper template - the
work product to which this Methods section belongs - and the subsequent
population of that template into the first full version of this paper were
performed by Claude Code Opus 4.7 Max with a 1M token context window using
the Max 5x plan \cite{claude-opus-47, claude-code}. The template-creation
pass produced the LaTeX skeleton, the style file, the bibliography, the
section instruction blocks, and this Methods section as a single multi-
commit pull request to the kevinkawchak/physical-ai-oncology-trials
repository \cite{main-repo}.

Additional color diagrams were also generated by Claude Code Opus 4.7 Max
(1M tokens) in a second future step that follows the template-population
pass. Paper text white space formatting and table column width
optimizations were performed by the author of the paper as a third future
step, applying the senior-author formatting rules captured in the global
formatting brief at the top of main.tex.

The current first prompt was written by the author of the paper, who also
created the color diagrams prompt that drives the second future step.
ChatGPT Thinking 5.5 with Extended Thinking and Deep Research, accessed
through ChatGPT Plus, was used to generate both the Background-A and
Background-B summaries that seed the bibliography and the Introduction
\cite{chatgpt-thinking}. Claude Sonnet 4.6 with Adaptive Thinking was used
to produce three chunks and a README file separately for both Background-A
and Background-B, yielding the eight Markdown source files at
new-trial/national-24-7-trial/Background-A/ and new-trial/national-24-7-trial/Background-B/
\cite{claude-sonnet-46}. Google Gemini AI Overview, accessed through
google.com search, was utilized for additional search assistance during
the source-gathering phase \cite{google-gemini-overview}.

\subsection{Repository Inputs}
\label{sec:methods-inputs}

The four simulations referenced throughout this paper are located at the
following repository paths and are documented in the Results section:

\begin{itemize}[leftmargin=*]
\item Simulation 1 - Continuous Real-Time Clinical Trial:
\path{new-trial/national-24-7-trial/hour-00/} through
\path{hour-55/}, plus
\path{new-trial/national-24-7-trial/extra-hours/hour-56/} through
\path{hour-83/}, with each hour folder containing four Markdown files and
three ASCII diagrams. See \cite{kawchak_2026_19176370} for the related
site documentation package.
\item Simulation 2 - Single-Patient 10-Stage Journey:
\path{patient-journey/stage_01_prescreening.py} through
\path{patient-journey/stage_10_closeout.py}, with the prior paper at
\path{patient-journey/paper/patient_journey_paper.tex}
\cite{kawchak_2026_19119939}.
\item Simulation 3 - 24-Hour Autonomous Sponsor:
\path{sponsor/final_paper/scripts/hourly/sponsor_hour_00.py} through
\path{sponsor_hour_23.py}, with 53 core agents under
\path{sponsor/final_paper/scripts/core_agents/} and 75 ASCII diagrams under
\path{sponsor/final_paper/scripts/diagrams/} \cite{kawchak_2026_19396256}.
\item Simulation 4 - 168-Hour 7-Day Extension:
\path{sponsor/final_paper/168_hours/day_01/} through
\path{day_07/}, with hourly scripts at
\path{day_NN/hourly/sponsor_hour_NNN.py} and the local-verification
instructions at
\path{sponsor/final_paper/168_hours/instructions/core_i5_6200u_4gb/}
\cite{kawchak_2026_19396256}.
\end{itemize}

The deep-research summaries that anchor the Introduction are located at
\path{new-trial/national-24-7-trial/Background-A/} (three Markdown chunks
plus a README, 17 BibTeX entries) and
\path{new-trial/national-24-7-trial/Background-B/} (three Markdown chunks
plus a README, 9 BibTeX entries; one entry, Yoo2025SCORPIO, is shared with
Background-A and is included once in the bibliography). The FDA April 28,
2026 source release is located at
\path{new-trial/national-24-7-trial/FDA-April-2026/FDA_RealTime_Clinical_Trials.md}
\cite{fda2026realtime}.

\subsection{Build and Verification}
\label{sec:methods-build}

The LaTeX paper compiles in Overleaf and locally with the standard four-
pass sequence:
\begin{verbatim}
pdflatex main.tex
bibtex main
pdflatex main.tex
pdflatex main.tex
\end{verbatim}

The bibliography uses the ieeetr style. Every reference includes a DOI
where one exists and a clickable URL through the \texttt{note} field. The
template ships an Overleaf-ready ZIP at
\path{new-trial/national-24-7-trial/paper/LaTeX_Source_Files.zip} that
contains \texttt{main.tex}, \texttt{new\_paper.sty}, \texttt{references.bib},
\texttt{orcid\_icon.png}, and the \texttt{sections/} directory.

\subsection{Continuous Integration Compatibility}
\label{sec:methods-ci}

The paper directory contains only LaTeX, Markdown, and PNG assets. No new
Python files are introduced by the template-creation pass, so the
repository's existing \texttt{ruff} configuration in \texttt{ruff.toml} and
the \texttt{lint-and-format} workflow at \path{.github/workflows/ci.yml}
across Python 3.10, 3.11, and 3.12 remain unaffected. The next Claude Code
4.7 Max pass that populates body sections must continue to add LaTeX
content only; if any helper Python script is later required, it must be
formatted with \texttt{ruff format} and pass \texttt{ruff check} under the
repository's 120-character line length and \texttt{[E, F, W]} rule
selection.
